% Generated by tools/edn_latex.py from reports/edn.md.
% Do not edit: change reports/edn.md and run `make edn`.
\ednentry{
  id = {EDN-41},
  title = {\texorpdfstring{\texttt{model\_\allowbreak{}version}}{model\_version} becomes the leading token of the normalised trim plus a frequency floor},
  date = {2026-09-30},
  milestone = {M2: Reproducibility (feature engineering; feeds M3: Quality Assurance)},
  topic = {Feature Encoding},
  participants = {@lukas2510},
  decision = {\texttt{model\_\allowbreak{}version} enters the extended feature set as the first \texttt{features.\allowbreak{}model\_\allowbreak{}version\_\allowbreak{}tokens} tokens of its normalised form -- case and accents folded, separators collapsed -- keeping only the levels above \texttt{features.\allowbreak{}model\_\allowbreak{}version\_\allowbreak{}min\_\allowbreak{}frequency} of the training rows and treating everything else as missing. Both numbers stay parameters so the experiment ladder can sweep them. The two finer variants measured below are rejected for now, and the justification in the code says what the token actually is rather than what it was assumed to be.},
  rationale = {\emph{Alternatives considered:}
\begin{itemize}[leftmargin=*, topsep=2pt]
\item \textbf{Option A (chosen): normalise, keep the leading token, then a frequency floor.}
\newline Pros: 279 levels covering 85.9\,\% of the training rows, from a column with 80,332 distinct raw values, at the cost of one regular expression. Both knobs are parameters, so the ladder measures whether the coarseness costs accuracy rather than anyone arguing about it. The level set is an artefact the API loads, so a request is normalised exactly as a training row was.
\newline Cons: coarse, and unevenly so. The leading token names the body style for Audi and Porsche (\texttt{avant}, \texttt{coupe}) and the engine letter for the German premium models that are most of the data: \texttt{d}, \texttt{911} and \texttt{e} are the three most frequent levels and 12.4\,\% of the training rows lead with an engine letter, which largely repeats \texttt{fuel\_\allowbreak{}category} while discarding the trim a buyer cares about (\texttt{M Sport}, \texttt{Touring}). Inside one make and model the merging is heavy -- Porsche 992 level \texttt{911} covers 743 distinct raw trims, Audi A6 \texttt{avant} 473, BMW 320 \texttt{d} 331 -- and 329 pairs of levels survive where one is a prefix of the other, so \texttt{(\textquotesingle{}20d\textquotesingle{}, \textquotesingle{}320d\textquotesingle{})} are separate levels for the same car. 14.1\,\% of the rows have no level at all.
\item \textbf{Option B: prefix the token with make and model, so a level is \texttt{BMW|320|d}.}
\newline Pros: reads as the obvious fix for the prefix-pair problem, and it does reduce those from 329 pairs to 90.
\newline Cons: measured, it does not fix the merging it was meant to fix and it costs a lot. The worst merge inside one make and model stays at exactly 743, because Porsche 992 level \texttt{911} is already inside one make and model, and the number of groups holding over 50 raw trims is unchanged at 203. Meanwhile each qualified level is rarer, so the same frequency floor drops far more of them: coverage falls from 85.9\,\% to 62.8\,\%, and the median merge rises from 3 to 47.
\item \textbf{Option C: keep two tokens when the first is an engine letter, one otherwise.}
\newline Pros: the only variant that touches the specific weakness. BMW 320's worst merge falls from 331 to 200, BMW 520's from 327 to 130, Mercedes C 220's from 74 to 35, and the groups over 50 raw trims fall from 203 to 179.
\newline Cons: not clearly better overall. It does nothing for the two largest merges (Porsche 992 \texttt{911} at 743, Audi A6 \texttt{avant} at 473, neither an engine letter), it costs 3 points of coverage (85.9\,\% to 82.9\,\%) and takes the level count from 279 to 318, and 1,797 training rows (2.9\,\%) lose their level entirely. It also hard-codes a per-make letter list into a rule that is otherwise a swept parameter, so it cannot be turned off from \texttt{params.\allowbreak{}yaml}.
\item \textbf{Option D: two tokens throughout, which is option A at \texttt{model\_\allowbreak{}version\_\allowbreak{}tokens: 2}.}
\newline Pros: the best merging of the four -- worst merge 358 rather than 743, only 64 groups over 50 trims -- and it needs no code at all, because it is a parameter change.
\newline Cons: coverage collapses to 44.0\,\%, so more than half the column is missing. Kept as the parameter value the ladder will try, not as the default.
\item \textbf{Option E: normalise only, with no token cut.}
\newline Pros: loses nothing.
\newline Cons: unusable as a level set. 78,738 distinct values survive normalisation over 105,405 listings, and a frequency floor on the full string keeps 7 levels covering 0.8\,\% of the rows.
\item \textbf{Option F: hash the trim into a fixed number of buckets.}
\newline Pros: bounded level count with no vocabulary to version.
\newline Cons: destroys the SHAP explainability FR-08 promises, because a bucket has no meaning to show a user.
\end{itemize}
\par
\emph{Why:} A is the cheapest thing that produces a usable level set, and the alternatives were measured rather than argued. B and C both looked better on the face of it and neither survived the measurement: B leaves the merging it targets untouched while costing 23 points of coverage, and C buys a real improvement on 12.4\,\% of the rows by making the other 87.6\,\% slightly worse and by hard-coding a rule the ladder cannot sweep. The decision that actually matters is that both knobs are parameters, so the question ``is the coarse trim good enough'' is answered by the experiment ladder with an MdAPE rather than by anyone's judgement here, and D is already the configuration it will compare against. What had to change regardless of the outcome is the justification: the code claimed the leading token ``is the one that names the trim'', which is true for Audi and Porsche body styles and false for the German premium engine codes that dominate the data, and a comment that claims a property the data does not have is worse than no comment.},
  ai = {Information seeking, Alternative generation, Alternative assessment, Recommendation},
  response = {Accepted with modifications},
  assessment = {The normalisation and the frequency floor were AI's own proposal in the \texttt{features} ticket, and the overstated justification was AI's wording. An adversarial review, also by AI, measured the collisions on the real column and showed the claim was false for most of the data; it proposed B and C as cheap improvements and asked for them to be implemented only if they measured better. They were then measured, and both were rejected on their own numbers -- which is the part worth recording, because the suggestion was plausible enough that it would have been accepted without the measurement. Lukas accepted keeping A with the corrected justification.},
  aievidence = {Claude Code session on 2026-09-30: an adversarial review of \href{https://github.com/mlops-2627q1-mds-upc/MLOPS_Recommenditos/pull/55}{PR \#55} reported that ``the docstring's claim that the leading token is the one that names the trim is true for Audi and Porsche body styles and false for the German premium engine codes, which are the bulk of the data'', and asked for the two variants to be evaluated and implemented ``only if it is clearly better on the real data and you can show the numbers''. The numbers are in \texttt{reports/\allowbreak{}analysis/\allowbreak{}extended\_\allowbreak{}features\_\allowbreak{}results.\allowbreak{}txt}.},
  evidence = {\href{https://github.com/mlops-2627q1-mds-upc/MLOPS_Recommenditos/blob/main/reports/analysis/extended_features.py}{\texttt{reports/\allowbreak{}analysis/\allowbreak{}extended\_\allowbreak{}features.\allowbreak{}py}} and its results file (2026-09-30); \href{https://github.com/mlops-2627q1-mds-upc/MLOPS_Recommenditos/blob/main/docs/docs/model-card.md}{model card} ``Feature space''; \href{https://github.com/mlops-2627q1-mds-upc/MLOPS_Recommenditos/blob/main/docs/docs/problem-spec.md}{problem specification} section 4; \hyperref[EDN-15]{EDN-15}; \href{https://github.com/mlops-2627q1-mds-upc/MLOPS_Recommenditos/issues/36}{issue \#36}; \href{https://github.com/mlops-2627q1-mds-upc/MLOPS_Recommenditos/pull/55}{PR \#55}.},
}
